\section{Historical Evidence and Scope}
\label{sec:historical-evidence}
Two archived campaigns explain the design's evolution. The original scalar
FFT study compares immediate and incremental execution at fixed requested
cadence. It locates residual preparation/reconstruction bursts and shows why
p99 can miss work occurring less than once per hundred sample calls.
Appendices~\ref{sec:method} and~\ref{sec:legacy-results} preserve its methods,
within-session resampling, and complete numerical and timing results.

The complete-pipeline prototype compares a serial one-hop schedule with
boundary bursts and four-hop overlap (Appendix~\ref{sec:pipeline-evaluation}).
It illustrates lower observed block maxima alongside higher aggregate cost;
the one-hop design also avoids the overlapped prototype's additional $3H$
samples of age. Some controls change storage and arithmetic as well as
placement. These historical results do not rank the current weighted analyzer
or optimized libraries. The confirmation study supplies that separate
evidence; neither historical nor current campaigns measure a running audio
device or concurrent display consumption.
